\section*{Bab \ref{ch:b1:precipitation} --- Pengendapan dan Kelarutan}

\begin{solution}{exo:b1:precipitation:1}
\ce{BaSO4(s) <=> Ba^2+ + SO4^2-}, $K_s = [\ce{Ba^2+}][\ce{SO4^2-}] =
10^{-9.97}$. \ce{CaF2(s) <=> Ca^2+ + 2F-}, $K_s = [\ce{Ca^2+}][\ce{F-}]^2 =
10^{-9.84}$. \ce{Ag2CrO4(s) <=> 2Ag+ + CrO4^2-}, $K_s =
[\ce{Ag+}]^2[\ce{CrO4^2-}] = 10^{-11.95}$. \ce{Fe(OH)3(s) <=> Fe^3+ + 3OH-},
$K_s = [\ce{Fe^3+}][\ce{OH-}]^3 = 10^{-38.55}$. Untuk $K_s = \num{2.0e-13}$,
$\mathrm{p}K_s = 12.70$.
\end{solution}

\begin{solution}{exo:b1:precipitation:2}
$s = \sqrt{K_s} = 10^{-12.27/2} = \qty{7.3e-7}{mol/L}$, yaitu
$\num{7.3e-7} \times 187.8 = \qty{1.4e-4}{g/L} = \qty{0.14}{mg/L}$.
\end{solution}

\begin{solution}{exo:b1:precipitation:3}
Sesudah pencampuran (volumenya menjadi dua kali), $[\ce{Ba^2+}] = \qty{1.0e-3}{mol/L}$ dan
$[\ce{SO4^2-}] = \qty{1.0e-4}{mol/L}$: $Q_s = \num{1.0e-7} > K_s =
\num{1.1e-10}$, barium sulfat mengendap. Pada kasus kedua
$[\ce{SO4^2-}] = \num{1.0e-7}$ dan $Q_s = \num{1.0e-10} < K_s$: tidak ada
endapan, walaupun nyaris.
\end{solution}

\begin{solution}{exo:b1:precipitation:4}
Di dalam air $s = 10^{-9.97/2} = \qty{1.0e-5}{mol/L}$. Di dalam sulfat
\qty{0.010}{mol/L} (\cref{prop:b1:precipitation:common-ion}), $s = 10^{-9.97}/0.010 =
\qty{1.1e-8}{mol/L}$, sekitar seribu kali lebih sedikit (faktor 970).
\end{solution}

\begin{solution}{exo:b1:precipitation:5}
Di dalam air $s = (10^{-9.84}/4)^{1/3} = \qty{3.3e-4}{mol/L}$, atau
\qty{26}{mg/L}. Di dalam kalsium klorida \qty{0.10}{mol/L},
$[\ce{Ca^2+}] \approx 0.10$ dan $[\ce{F-}] = 2s$, sehingga
$s = \frac12\sqrt{10^{-9.84}/0.10} = \qty{1.9e-5}{mol/L}$. Pendekatan
$s \ll 0.10$ berlaku dengan selisih empat orde besar.
\end{solution}

\begin{solution}{exo:b1:precipitation:6}
Awal pengendapan: $\mathrm{pH} = 14.00 - \frac12(11.25 + \log 0.050) = 14.00 - 4.97 =
9.03$. Pada \qty{99}{\percent}, $[\ce{Mg^2+}] = \num{5.0e-4}$ dan
$\mathrm{pH} = 14.00 - \frac12(11.25 - 3.30) = 10.03$. Pada sumbu pH,
\ce{Mg(OH)2} ada di atas 9.03; di bawahnya, magnesium seluruhnya berupa
\ce{Mg^2+} pada \qty{0.050}{mol/L}.
\end{solution}

\begin{solution}{exo:b1:precipitation:7}
Perak bromida muncul lebih dahulu, pada $[\ce{Ag+}] = 10^{-12.27 + 2} =
10^{-10.27}$, perak klorida pada $10^{-7.75}$. Ketika perak klorida
muncul, $[\ce{Br-}] = 10^{-12.27}/10^{-7.75} = 10^{-4.52} =
\qty{3.0e-5}{mol/L}$, \qty{0.30}{\percent} dari bromida. Dengan kriteria
yang biasa, \qty{0.1}{\percent}, pemisahan itu tidak dapat diterima:
kedua $\mathrm{p}K_s$ hanya berselisih 2.5.
\end{solution}

\begin{solution}{exo:b1:precipitation:8}
Perak kromat terbentuk ketika $[\ce{Ag+}]^2 \times \num{5.0e-3} = 10^{-11.95}$,
$[\ce{Ag+}] = \qty{1.5e-5}{mol/L}$. Maka $[\ce{Cl-}] = 10^{-9.75}/\num{1.5e-5}
= \qty{1.2e-5}{mol/L}$, bagian \num{2.4e-4} (\qty{0.024}{\percent})
klorida: warna merah muncul ketika klorida praktis seluruhnya telah
mengendap.
\end{solution}

\begin{solution}{exo:b1:precipitation:9}
$[\ce{CO3^2-}] = [\ce{HCO3-}]K_{a2}/[\ce{H3O+}] = [\ce{HCO3-}]\,10^{\mathrm{pH}
- \mathrm{p}K_{a2}}$, sehingga $K_s = [\ce{Ca^2+}][\ce{HCO3-}]\,10^{\mathrm{pH} -
\mathrm{p}K_{a2}}$, yaitu hubungan yang diminta. Dengan $[\ce{Ca^2+}] =
[\ce{HCO3-}] = s$: $s^2 = 10^{-8.30 + 10.33 - 8.30} = 10^{-6.27}$,
$s = \qty{7.3e-4}{mol/L}$ (\qty{73}{mg/L} \ce{CaCO3}), dibandingkan dengan
$10^{-4.15} = \qty{7.1e-5}{mol/L}$ seandainya ion karbonat bukan basa.
\end{solution}

\begin{solution}{exo:b1:precipitation:10}
1.~Dengan menambahkan \ce{CO2(g) <=> CO2(aq)}: $K' = K K_H = 10^{-4.34 - 1.47} =
10^{-5.81}$.
2.~$[\ce{Ca^2+}] = s$ dan $[\ce{HCO3-}] = 2s$, sehingga $4s^3 = K'p$. Pada
\qty{0.010}{bar}: $s = (10^{-5.81} \times 0.010/4)^{1/3} =
\qty{1.6e-3}{mol/L}$, \qty{157}{mg/L}. Di bawah udara luar:
$s = \qty{5.5e-4}{mol/L}$, \qty{55}{mg/L}. (Pemeriksaan: di bawah udara
tanah $[\ce{CO2(aq)}] = 10^{-3.47}$ dan $\mathrm{pH} = 6.37 + \log(3.14\times
10^{-3}/3.39 \times 10^{-4}) = 7.34$: hidrogenkarbonat memang
predominan.)
3.~Air yang jenuh kalsit di bawah tekanan karbon dioksida tanah yang
tinggi mencapai langit-langit gua, tempat tekanannya lebih rendah: air
itu kehilangan karbon dioksida, $Q > K'$, dan kalsit mengendap di tempat
tetesan itu bergantung, lingkaran demi lingkaran.
% ledger: air:CO2, dfg:CO2_g, dfg:CO2_ao (log KH computed in figdata/bachelor-1/aqueous-constants.py)
\end{solution}

\begin{solution}{exo:b1:precipitation:11}
1.~Awal pengendapan pada $\mathrm{pH} = 14.00 - \frac12(16.38 - 2.00) = 6.81$.
Pelarutan kembali yang sempurna terjadi ketika
$[\ce{Zn(OH)4^2-}] = 10^{-2} = 10^{-1.93}[\ce{OH-}]^2$,
$[\ce{OH-}] = 10^{-0.035}$, $\mathrm{pH} = 13.97$. Hidroksida itu ada
antara pH 6.81 dan 13.97.
2.~$\mathrm{pH} = 14.00 - 14.45/4 = 10.39$, $s_{\min} = 2 \times
10^{-(16.38 + 1.93)/2} = \qty{1.4e-9}{mol/L}$ (dalam model dua spesies
ini).
3.~Garis berkemiringan $-2$ dari $(6.81, -2)$ turun ke minimum di dekat
$(10.39, -8.9)$, lalu kemiringan $+2$ naik ke $(13.97, -2)$; kurva yang
sebenarnya dibulatkan oleh spesies-spesies hidrokso lainnya
(gambar pada \cref{sec:b1:precipitation:amphoteric}).
\end{solution}

\begin{solution}{exo:b1:precipitation:12}
1.~$\ce{Al(OH)3(s) + OH- <=> Al(OH)4-}$, $K = 10^{-1.23}$; semua aluminium
larut ketika $10^{-3} \leqslant K[\ce{OH-}]$, yaitu
$[\ce{OH-}] \geqslant 10^{-1.77}$, $\mathrm{pH} \geqslant 12.23$.
2.~$[\ce{Fe^3+}] = 10^{-38.55}/10^{-3 \times 1.77} = 10^{-33.2}$: besi
tetap seluruhnya sebagai \ce{Fe(OH)3}; penyaringan memisahkan besi
(padatan) dari aluminium (filtrat).
3.~Magnesium hidroksida tidak amfoter: dalam hidroksida berlebih ia
tetap padat, bersama besi(III) hidroksida, dan tidak ada yang terpisah.
\end{solution}

\begin{solution}{pb:b1:precipitation:1}
\textbf{1.} \ce{Fe(OH)3(s) <=> Fe^3+ + 3OH-}, $K_s = [\ce{Fe^3+}][\ce{OH-}]^3$;
\ce{Zn(OH)2(s) <=> Zn^2+ + 2OH-}, $K_s = [\ce{Zn^2+}][\ce{OH-}]^2$;
\ce{Mg(OH)2(s) <=> Mg^2+ + 2OH-}, serupa.
\textbf{2.} Tetapan-tetapannya tidak berbentuk sama (pangkat 3 atau 2
dari $[\ce{OH-}]$) dan ketiga ion itu berada pada konsentrasi yang
berbeda; hanya pH awal pengendapan yang dapat dibandingkan.
\textbf{3.} Pada awal pengendapan $c[\ce{OH-}]^n = K_s$, sehingga
$n\log[\ce{OH-}] =
-\mathrm{p}K_s - \log c$ dan $\mathrm{pH} = \mathrm{p}K_e +
\log[\ce{OH-}] = \mathrm{p}K_e - \frac1n(\mathrm{p}K_s + \log c)$.
\textbf{4.} $14.00 - \frac13(38.55 - 3.00) = 2.15$.
\textbf{5.} $14.00 - \frac12(16.38 - 2.30) = 6.96$.
\textbf{6.} $14.00 - \frac12(11.25 - 1.70) = 9.22$.
\textbf{7.} Besi(III) (2.15), lalu seng (6.96), lalu magnesium (9.22). Pada
pH 2.0, $Q_s = 10^{-3} \times 10^{-36} = 10^{-39} < 10^{-38.55}$ untuk
besi, dan yang lain lebih jauh dari mengendap: tidak ada yang padat.
\textbf{8.} \qty{1.0e-6}{mol/L}.
\textbf{9.} $14.00 - \frac13(38.55 - 6.00) = 3.15$.
\textbf{10.} $\log[\ce{Fe^3+}] = -38.55 + 3(14.00 - \mathrm{pH}) = 3.45 -
3\,\mathrm{pH}$: kemiringan $-3$.
\textbf{11.} Sampai awal pengendapan seng hidroksida, pH 6.96.
\textbf{12.} Dari 3.15 sampai 6.96.
\textbf{13.} \ce{Fe^3+ + OH- <=> FeOH^2+}, $\beta_1 =
[\ce{FeOH^2+}]/([\ce{Fe^3+}][\ce{OH-}]) = 10^{11.82}$.
\textbf{14.} Keduanya sama ketika $[\ce{OH-}] = 10^{-11.82}$, pada pH 2.18:
\ce{Fe^3+} predominan di bawahnya, \ce{FeOH^2+} di atasnya.
\textbf{15.} $10^{11.82 + 3.15 - 14.00} = 10^{0.97} = 9.3$: pada pH 3.15
besi terlarut 10.3 kali \ce{Fe^3+} bebas, \qty{1.0e-5}{mol/L}: hanya
\qty{99.0}{\percent} yang tersingkirkan. Perkiraan pertama itu tidak aman.
\textbf{16.} $[\ce{Fe}]_{\mathrm{terlarut}} = 10^{3.45 - 3\,\mathrm{pH}}\,
(1 + 10^{\mathrm{pH} - 2.18})$.
\textbf{17.} $\log[\ce{FeOH^2+}] = 11.82 + \log[\ce{Fe^3+}] + \mathrm{pH} -
14.00 = 1.27 - 2\,\mathrm{pH}$: kemiringan $-2$.
\textbf{18.} $1.27 - 2\,\mathrm{pH} \approx -6.00$ menghasilkan 3.64;
dengan suku \ce{Fe^3+} disertakan, $10^{3.45 - 10.92}(1 + 10^{1.46}) =
\num{1.0e-6}$ pada $\mathrm{pH} = 3.64$.
\textbf{19.} $[\ce{Zn^2+}] = \num{5.0e-5}$: $14.00 - \frac12(16.38 - 4.30) =
7.96$. Magnesium hidroksida baru mulai pada 9.22: tidak.
\textbf{20.} \ce{Zn(OH)2(s) + 2OH- <=> Zn(OH)4^2-}, $K = \beta_4 K_s =
10^{14.45 - 16.38} = 10^{-1.93}$.
\textbf{21.} $[\ce{OH-}]^2 = \num{5.0e-3}/10^{-1.93}$, $\mathrm{pH} =
13.81$. Operator tidak boleh menambahkan basa terlalu banyak: di atas
pH 9.2 magnesium mengendap bersama seng, dan jauh di atasnya, seng
kembali larut.
\textbf{22.} $\num{1.0e-3} \times 1000 \times 0.999 \times 106.8 =
\qty{107}{g}$ per meter kubik.
\textbf{23.} $\num{5.0e-3} \times 1000 \times 0.99 \times 99.4 =
\qty{492}{g}$ per meter kubik.
\textbf{24.} Jika pH dinaikkan sekaligus di atas 6.96, kedua hidroksida
mengendap dalam satu lumpur: seng hilang dalam limbah besi, dan limbahnya
tercemar seng. Berhenti di antara kedua ambang menyingkirkan besi saja.
\textbf{25.} \textbf{Dari pH 3.64 sampai pH 6.96}: di bawah 3.64 lebih
dari \qty{0.1}{\percent} besi masih terlarut, sebagian besar sebagai
\ce{FeOH^2+}; di atas 6.96 seng hidroksida terbentuk. Mengabaikan
\ce{FeOH^2+} akan menempatkan batas bawahnya pada 3.15.
\end{solution}
